% GENERATED FILE. Edit canonical lessons, then run npm run generate:books.
% canonical-source-hash: fnv1a64:22235a5e94ac48b4
% canonical-lessons: TE-C81-first, TE-C81-second, TE-C81-third, TE-C81-fourth, TE-C81-fifth

\chapter{Counting in Order}
\label{ch:te-ordinals}
\hlchaptermodality{81}

\begin{chapteropening}
\textbf{By the end of this chapter:} \emph{I can say first to fifth in Telugu, and make an ordinal from any number by dropping its final vowel and adding -va.}

The five said in order, with only modati learned as a word and the other four made by the ending.
\end{chapteropening}

\section[modaṭi]{\te{మొదటి} (modaṭi) — first}

\label{lesson:TE-C81-first}



\begin{quote}
A new kind of counting starts here, so the recalls come first, at the four growing distances this book measures.

\begin{itemize}
\raggedright
  \item \emph{Recall:} say \emph{śubhākāṅkṣalu} — \textbf{R1}, one lesson back, before it decays
  \item \emph{Recall:} say \emph{atanu} — \textbf{R2}, the first real retrieval, five to fifteen lessons back
  \item \emph{Recall:} say \emph{goḍḍali} — \textbf{R3}, durable, twenty to sixty lessons back
  \item \emph{Recall:} say \emph{atta} — \textbf{R3}, a second word from the same distance
  \item \emph{Recall:} say \emph{padakoṇḍu and iravai} — \textbf{R4}, recognition at distance, eighty lessons back or more
  \item \emph{Recall:} say \emph{kukka and pilli} — \textbf{R4}, a second word from the same distance
\end{itemize}
\end{quote}

\subsection*{You'll want to know: \te{మొదటి}}
\textbf{\te{మొదటి}} (\emph{modaṭi}) — \textquotedblleft{}\textbf{first}\textquotedblright{}.

\emph{\te{ఒకటి}, \te{రెండు}, \te{మూడు}} answer \textbf{how many}. This word answers \textbf{which one in the line}, and Telugu keeps a second set of words for that job.

It goes straight in front of what it counts: \textbf{\te{మొదటి} \te{రోజు}} (\emph{modaṭi rōju}), the first day; \textbf{\te{మొదటి} \te{పుస్తకం}} (\emph{modaṭi pustakaṁ}), the first book. Nothing changes on the noun.

\begin{cousinweb}
There is no \textbf{\te{ఒకటి}} in \textbf{\te{మొదటి}}, and there was never going to be. It is built on \textbf{\te{మొదలు}} (\emph{modalu}), \textquotedblleft{}\textbf{a beginning}\textquotedblright{} — so \emph{modaṭi} means \textquotedblleft{}of the beginning\textquotedblright{}, not \textquotedblleft{}one-th\textquotedblright{}. You will also meet the bare \textbf{\te{మొదట}} (\emph{modaṭa}), \textquotedblleft{}at first\textquotedblright{}.

That is worth noticing because \emph{\te{ఒకటి}} was already the maverick of the cardinals: Tamil says \emph{oṉṟu}, Kannada \emph{ondu}, Malayalam \emph{onnu}, and Telugu alone says \emph{okaṭi}. Now the number \textbf{one} turns out to be doubly on its own — the sisters do not share its cardinal, and Telugu does not build its ordinal on it either.

From two upward it is a different story, and the next lesson is where the rule starts.
\end{cousinweb}

\subsection*{Your turn}
Say these aloud:

\begin{itemize}
\raggedright
  \item \emph{modaṭi}
  \item \emph{modaṭi rōju} — the first day
  \item \emph{okaṭi} is how many, \emph{modaṭi} is which one
\end{itemize}

\begin{tcolorbox}[breakable,colback=teal!4,colframe=teal!35!black,title={Before you move on}]
What does \te{మొదటి} mean? (\textbf{First} — which one in the line.) What word is it built on? (\textbf{\te{మొదలు}}, a beginning — \textbf{not} \te{ఒకటి}.) Say \textquotedblleft{}the first day\textquotedblright{}. (\textbf{\te{మొదటి} \te{రోజు}}.)
\end{tcolorbox}
\section[reṇḍava]{\te{రెండవ} (reṇḍava) — second, and the swap behind all the rest}

\label{lesson:TE-C81-second}



\begin{quote}
Recalls before the second ordinal, at the four measured distances.

\begin{itemize}
\raggedright
  \item \emph{Recall:} say \emph{modaṭi} — \textbf{R1}, one lesson back, before it decays
  \item \emph{Recall:} say \emph{nannu} — \textbf{R2}, the first real retrieval, five to fifteen lessons back
  \item \emph{Recall:} say \emph{rampaṁ} — \textbf{R3}, durable, twenty to sixty lessons back
  \item \emph{Recall:} say \emph{bāva} — \textbf{R3}, a second word from the same distance
  \item \emph{Recall:} say \emph{the letter e} — \textbf{R4}, recognition at distance, eighty lessons back or more
  \item \emph{Recall:} say \emph{pacca and pasupu} — \textbf{R4}, a second word from the same distance
\end{itemize}
\end{quote}

\subsection*{You'll want to know: \te{రెండవ}}
\textbf{\te{రెండవ}} (\emph{reṇḍava}) — \textquotedblleft{}\textbf{second}\textquotedblright{}.

\textbf{\te{రెండవ} \te{రోజు}} (\emph{reṇḍava rōju}) — the second day.

\begin{grammarlens}[title={Grammar Lens: cardinal, minus its last vowel, plus -\te{వ}}]
Say them together and you can hear the join: \textbf{\te{రెండు}} $\to$ \textbf{\te{రెండవ}}.

Take the number you already have, take its final vowel off, and put \textbf{-\te{వ}} (\emph{-va}) on the end. That is the whole of it, and it holds for every number from two upward. Nine ordinals for the price of one ending.

Telugu is doing here what the last lesson said it would not do for \textbf{one}: \emph{modaṭi} had no number to build on, and every ordinal from here does.

In speech you will also hear \textbf{\te{రెండో}} (\emph{reṇḍō}), the same word ending in a long \textbf{ō} instead of \textbf{-\te{వ}}. It is the everyday spoken form; \textbf{-\te{వ}} is what you will see written.
\end{grammarlens}

\subsection*{Your turn}
Say these aloud:

\begin{itemize}
\raggedright
  \item \emph{reṇḍava}
  \item the swap — \emph{reṇḍu} $\to$ \emph{reṇḍava}
  \item \emph{reṇḍava rōju} — the second day
\end{itemize}

\begin{tcolorbox}[breakable,colback=teal!4,colframe=teal!35!black,title={Before you move on}]
How is an ordinal made from two upward? (\textbf{Take the number, drop its final vowel, add -\te{వ}}.) Say the second day. (\textbf{\te{రెండవ} \te{రోజు}}.) And what will you hear instead in speech? (\textbf{\te{రెండో}}.)
\end{tcolorbox}
\section[mūḍava]{\te{మూడవ} (mūḍava) — third}

\label{lesson:TE-C81-third}



\begin{quote}
Recalls first, then the third.

\begin{itemize}
\raggedright
  \item \emph{Recall:} say \emph{reṇḍava} — \textbf{R1}, one lesson back, before it decays
  \item \emph{Recall:} say \emph{gāru} — \textbf{R2}, the first real retrieval, five to fifteen lessons back
  \item \emph{Recall:} say \emph{pāra} — \textbf{R3}, durable, twenty to sixty lessons back
  \item \emph{Recall:} say \emph{maradalu} — \textbf{R3}, a second word from the same distance
  \item \emph{Recall:} say \emph{the letter ṇa} — \textbf{R4}, recognition at distance, eighty lessons back or more
  \item \emph{Recall:} say \emph{the letter ja} — \textbf{R4}, a second word from the same distance
\end{itemize}
\end{quote}

\subsection*{You'll want to know: \te{మూడవ}}
\textbf{\te{మూడవ}} (\emph{mūḍava}) — \textquotedblleft{}\textbf{third}\textquotedblright{}.

\textbf{\te{మూడు}} $\to$ \textbf{\te{మూడవ}}. The long \textbf{ū} stays long; only the ending changes.

\textbf{\te{మూడవ} \te{వారం}} (\emph{mūḍava vāraṁ}) — the third week.

Two numbers in a row have now taken the same ending, which is what a rule looks like from inside: not a statement, but the second time you were not surprised. The spoken form is \textbf{\te{మూడో}} (\emph{mūḍō}), exactly as \emph{reṇḍō} was.

\subsection*{Your turn}
Say these aloud:

\begin{itemize}
\raggedright
  \item \emph{mūḍava}
  \item the swap — \emph{mūḍu} $\to$ \emph{mūḍava}
  \item \emph{mūḍava vāraṁ} — the third week
\end{itemize}

\begin{tcolorbox}[breakable,colback=teal!4,colframe=teal!35!black,title={Before you move on}]
Say the third. (\textbf{\te{మూడవ}}.) What happened to the long \emph{ū}? (\textbf{Nothing} — only the ending changed.) And the spoken form? (\textbf{\te{మూడో}}.)
\end{tcolorbox}
\section[nālugava]{\te{నాలుగవ} (nālugava) — fourth}

\label{lesson:TE-C81-fourth}



\begin{quote}
Recalls first, then the fourth.

\begin{itemize}
\raggedright
  \item \emph{Recall:} say \emph{mūḍava} — \textbf{R1}, one lesson back, before it decays
  \item \emph{Recall:} say \emph{nāku kāvāli} — \textbf{R2}, the first real retrieval, five to fifteen lessons back
  \item \emph{Recall:} say \emph{sutti} — \textbf{R3}, durable, twenty to sixty lessons back
  \item \emph{Recall:} say \emph{alluḍu} — \textbf{R3}, a second word from the same distance
  \item \emph{Recall:} say \emph{śubha rātri} — \textbf{R4}, recognition at distance, eighty lessons back or more
  \item \emph{Recall:} say \emph{śubha sāyantram} — \textbf{R4}, a second word from the same distance
\end{itemize}
\end{quote}

\subsection*{You'll want to know: \te{నాలుగవ}}
\textbf{\te{నాలుగవ}} (\emph{nālugava}) — \textquotedblleft{}\textbf{fourth}\textquotedblright{}.

\textbf{\te{నాలుగు}} $\to$ \textbf{\te{నాలుగవ}}. Four syllables, and the swap is the same one: the final \textbf{u} goes, \textbf{-\te{వ}} arrives.

\textbf{\te{నాలుగవ} \te{సంవత్సరం}} (\emph{nālugava saṁvatsaraṁ}) — the fourth year.

This is the longest of the ten and it is no harder than the shortest, which is the practical point of a regular ending. Once you can count, you can order.

\subsection*{Your turn}
Say these aloud:

\begin{itemize}
\raggedright
  \item \emph{nālugava}
  \item the swap — \emph{nālugu} $\to$ \emph{nālugava}
  \item \emph{nālugava saṁvatsaraṁ} — the fourth year
\end{itemize}

\begin{tcolorbox}[breakable,colback=teal!4,colframe=teal!35!black,title={Before you move on}]
Say the fourth. (\textbf{\te{నాలుగవ}}.) What is the ending? (\textbf{-\te{వ}}, after the final vowel goes.) And say \textquotedblleft{}the fourth year\textquotedblright{}. (\textbf{\te{నాలుగవ} \te{సంవత్సరం}}.)
\end{tcolorbox}
\section[aidava]{\te{ఐదవ} (aidava) — fifth}

\label{lesson:TE-C81-fifth}



\begin{quote}
Recalls first, then the fifth, which closes the chapter.

\begin{itemize}
\raggedright
  \item \emph{Recall:} say \emph{nālugava} — \textbf{R1}, one lesson back, before it decays
  \item \emph{Recall:} say \emph{ani} — \textbf{R2}, the first real retrieval, five to fifteen lessons back
  \item \emph{Recall:} say \emph{karra} — \textbf{R3}, durable, twenty to sixty lessons back
  \item \emph{Recall:} say \emph{aṅgaḍi} — \textbf{R3}, a second word from the same distance
  \item \emph{Recall:} say \emph{śubha madhyāhnam} — \textbf{R4}, recognition at distance, eighty lessons back or more
  \item \emph{Recall:} say \emph{when to use śubha madhyāhnam} — \textbf{R4}, a second word from the same distance
\end{itemize}
\end{quote}

\subsection*{You'll want to know: \te{ఐదవ}}
\textbf{\te{ఐదవ}} (\emph{aidava}) — \textquotedblleft{}\textbf{fifth}\textquotedblright{}.

\textbf{\te{ఐదు}} $\to$ \textbf{\te{ఐదవ}}. The opening \textbf{\te{ఐ}} is untouched — it is a single letter in Telugu, not two — and only the final \textbf{u} is traded for \textbf{-\te{వ}}.

\textbf{\te{ఐదవ} \te{రోజు}} (\emph{aidava rōju}) — the fifth day.

That is one to five in order. Say \emph{modaṭi, reṇḍava, mūḍava, nālugava, aidava}, and notice that only the first one had to be learned as a word.

\subsection*{Your turn}
Say these aloud:

\begin{itemize}
\raggedright
  \item \emph{aidava}
  \item the five in order — \emph{modaṭi, reṇḍava, mūḍava, nālugava, aidava}
  \item \emph{aidava rōju} — the fifth day
\end{itemize}

\begin{tcolorbox}[breakable,colback=teal!4,colframe=teal!35!black,title={Before you move on}]
Say the five in order. (\emph{\textbf{Modaṭi, reṇḍava, mūḍava, nālugava, aidava}}.) How many of them had to be learned as separate words? (\textbf{One} — \emph{modaṭi}.) And what did the other four need? (\textbf{The number, and -\te{వ}}.)
\end{tcolorbox}
